\chapter{Where the Weil class is known to be algebraic}\label{ch:known}

This chapter collects the cases in which the Weil classes are known to be
algebraic. They fall into three groups. At special members of a Weil family,
where the endomorphism algebra is larger than $K$, the Weil classes can be
products of divisor classes, and we prove this in an explicit example. For
the field $\QQ(\sqrt{-3})$, and for some families for $\QQ(\sqrt{-1})$,
curves with an automorphism give cycles on Prym varieties, by work of
Schoen and Koike. And in dimension four, Markman's secant sheaves settle
every family. The last section reports, separately and conditionally, what
recent unrefereed preprints announce.

We use the notation of Chapter~\ref{ch:weil}. In particular $\mathrm{Sh}$
is the base of a Weil family of dimension $2n$, $\Sigma\subseteq\mathrm{Sh}$
is the set of members on which the Weil classes are algebraic, and by
Corollary~\ref{cor:oneclass} it does not matter whether we ask for one
nonzero Weil class or for all of them.

\section{A member where divisors suffice}\label{sec:cmpoint}

Let $E$ be an elliptic curve with complex multiplication by $K$, so that
$\End^{0}(E)=K$. Let $z$ be a linear coordinate on its universal cover, and
normalise the action so that $\tau\in K$ acts on the holomorphic form $dz$
by $\tau$, through the fixed embedding $K\subset\CC$. Then $\tau$ acts on
$d\bar z$ by $\bar\tau$.

On $A=E^{2n}$ let $\tau\in K$ act by $\tau$ on the first $n$ factors and by
$\bar\tau$ on the last $n$. Give $A$ the product polarisation. The Rosati
involution of $E$ is complex conjugation on $K$, so the Rosati involution of
$A$ restricts to complex conjugation on this copy of $K$.

\begin{proposition}\label{prop:cmexample}
With this action $A$ is of Weil type for $K$, and its Weil classes are
polynomials in divisor classes. In particular they are algebraic.
\end{proposition}

\begin{proof}
Write $z_{1},\dots,z_{2n}$ for the coordinates on the factors. The element
$\tau$ acts by $\tau$ on $dz_{i}$ for $i\le n$, and on $d\bar z_{i}$ for
$i>n$, because on those factors it acts through $\bar\tau$, whose conjugate
is $\tau$. So
\[
  V_{+}=\langle dz_{1},\dots,dz_{n},d\bar z_{n+1},\dots,d\bar z_{2n}\rangle,
\]
which has $n$ holomorphic and $n$ antiholomorphic directions, and $A$ is of
Weil type. Up to sign,
\[
  \textstyle\bigwedge^{2n}V_{+}=\prod_{i=1}^{n}\bigl(dz_{i}\wedge d\bar z_{n+i}\bigr).
\]
Each factor $dz_{i}\wedge d\bar z_{n+i}$ lies in the $(1,1)$ part of
$H^{1}(E_{i})\otimes H^{1}(E_{n+i})$, the Künneth component of
$H^{2}(E\times E)$ that corresponds to $\Hom(E,E)$. That $(1,1)$ part has
complex dimension two, spanned by $dz_{i}\wedge d\bar z_{n+i}$ and
$d\bar z_{i}\wedge dz_{n+i}$. On the other hand, the classes of the graphs of
the endomorphisms $1$ and $\sqrt{-d}$ of $E$, minus their components in
$H^{0}\otimes H^{2}$ and $H^{2}\otimes H^{0}$, give two linearly independent
rational classes in that component. So the whole $(1,1)$ part is spanned
over $\CC$ by divisor classes, and $\bigwedge^{2n}V_{+}$ lies in the
complex span of products of $n$ divisor classes. The same holds for
$\bigwedge^{2n}V_{-}$, its complex conjugate. So $\HW\otimes\CC$ lies in
$\mathcal D^{n}(A)\otimes\CC$. Both $\HW$ and $\mathcal D^{n}(A)$ are
rational subspaces, so $\HW\subseteq\mathcal D^{n}(A)$.
\end{proof}

The example has one discriminant, determined by the product form. A
construction with quaternion algebras gives such a member in every Weil
family, of every dimension and every discriminant, with the Weil classes
again polynomials in divisor classes \cite{BhH8}. We record the consequence
and prove the part we need.

\begin{proposition}\label{prop:densebook}
Suppose $\Sigma$ contains one point of $\mathrm{Sh}$. Then $\Sigma$ is dense
in $\mathrm{Sh}$.
\end{proposition}

\begin{proof}
Let $s_{0}\in\Sigma$ lift to $\tilde s_{0}\in D$. An element $g$ of
$\SU(V,H)(\QQ)$ is a $K$-linear automorphism of $V$ preserving $E$. It is an
isomorphism of rational Hodge structures from the structure at
$\tilde s_{0}$ to the structure at $g\tilde s_{0}$, so a positive integer
multiple of it is an isogeny between the corresponding abelian varieties.
The isogeny commutes with $K$, so it carries Weil classes to Weil classes,
and it carries cycles to cycles. So the image of $g\tilde s_{0}$ lies in
$\Sigma$. It remains to see that the orbit $\SU(V,H)(\QQ)\cdot\tilde s_{0}$
is dense in $D$. The real group $\SU(V,H)(\RR)=\SU(n,n)$ acts transitively
on $D$, so it suffices that $\SU(V,H)(\QQ)$ is dense in $\SU(V,H)(\RR)$. This
is real approximation for the connected group $\SU(V,H)$
\cite[Chapter~7]{PR94}.
\end{proof}

So the algebraic locus is never small in the topological sense. What is not
known is whether it has interior. Every point that this proof produces is
isogenous to $s_{0}$, so it has the same Hodge group as $s_{0}$, and
in the example above that group is a torus. Chapter~\ref{ch:remaining}
explains why this is the heart of the problem.

\section{Lowering the dimension}

The next statement, due to Schoen, lets a result in one dimension imply a
result one step lower.

\begin{proposition}[Schoen \cite{Sch98}]\label{prop:lower}
Let $A$ be of Weil type for $K$, of dimension $2n$ and discriminant
$\delta$, and let $B$ be an abelian surface of Weil type for $K$ with the
same discriminant. Then $A\times B$ is of Weil type for $K$, of dimension
$2n+2$, and its hermitian form is split, that is, it has a totally
isotropic $K$-subspace of dimension $n+1$. If the Weil classes of $A\times B$
are algebraic, so are those of $A$.
\end{proposition}

The first part is a computation with hermitian forms. The form of
$A\times B$ is the orthogonal sum of the two forms. With the discriminants
normalised as in \cite{BhH8}, matching discriminants make the sum have the
discriminant of a sum of $n+1$ hyperbolic planes. By Landherr's theorem, a
hermitian form over $K$ is determined by its rank, its discriminant and its
signature, so the sum is hyperbolic. For the
second part, the Weil plane of $A\times B$ is the $K$-tensor product of
the Weil planes of $A$ and $B$. The Weil classes of $B$ lie in $H^{2}$, where
they are divisor classes, and intersecting with a suitable divisor class
pulled back from $B$ and pushing forward to $A$ returns a nonzero Weil class
of $A$. A complete proof in coordinates is in \cite{BhH8}.

The consequence is that the \emph{split} families in dimension $2n+2$
control all families in dimension $2n$. This is how Markman passes from
sixfolds to fourfolds.

\section{Curves with an automorphism}\label{sec:schoen}

Schoen found the first cycles representing Weil classes at members with
full Hodge group. Let $\widetilde C\to C$ be an \'etale cyclic cover of degree
three of a smooth projective curve of genus $g\ge2$, with covering
automorphism $\rho$. The Prym variety $P$ is the abelian subvariety of
$J(\widetilde C)$ on whose first cohomology $\rho$ has no fixed vector.
Then $K=\QQ(\sqrt{-3})$ acts on $P$ through $\rho$. By the
Chevalley--Weil formula, the two nontrivial eigenspaces of $\rho$ on
holomorphic differentials of $\widetilde C$ both have dimension $g-1$, so $P$
is of Weil type with $n=g-1$. Its hermitian form is split, by a computation
of Looijenga with the homology of the cover \cite{Loo97}.

\begin{theorem}[Schoen \cite{Sch88}]\label{thm:schoen}
The Weil classes of these Prym varieties are algebraic.
\end{theorem}

Schoen's cycles come from the canonical system of the base curve. On the
symmetric power $C^{(2g-2)}$, the canonical system $|K_{C}|\cong\PP^{g-1}$ is
the fibre of the Abel--Jacobi map over the canonical class. The cover
$\widetilde C\to C$ defines an \'etale cover of $C^{(2g-2)}$ of degree three,
and over $|K_{C}|$, which is simply connected, it splits into three copies of
$\PP^{g-1}$. The differences of the copies, weighted by the two characters of
order three, are cycles on a quotient of $\widetilde C^{\,2g-2}$ whose classes
span the Weil classes there, and a correspondence through $J(\widetilde C)$
carries them to $P$ \cite[Theorem~2.0 with $r=0$, Corollary~3.1]{Sch88}. Patel and Zhang have recently given a second proof using unramified class
field theory (preprint) \cite{PZ25}. Looijenga showed that the monodromy of
the family of these Prym varieties is an arithmetic subgroup of the
unitary group \cite{Loo97}. By Borel density and Andr\'e's theorem, as in
Proposition~\ref{prop:generalsu}, a very general Prym of this kind has full
Hodge group $\SU(V,H)$. So Schoen's cycles live at members with full Hodge
group. We follow the account in \cite{BhH8}.

The family of Pryms has dimension $3g-3=3n$, the dimension of the moduli of
curves of genus $g$. The Weil family has dimension $n^{2}$. For $n=2$ Schoen
showed that the Pryms fill the Weil family of fourfolds of their polarization
type, and proved the conjecture there \cite[Theorem~3.2]{Sch88}; his addendum
points out that the theorem should have carried this hypothesis on the
polarization \cite{Sch98}. For $n=3$ the addendum shows, with a calculation of
Faber, that the Pryms fill a family of sixfolds, and then descends through
Proposition~\ref{prop:lower} to every Weil fourfold for $\QQ(\sqrt{-3})$,
whatever its polarization invariant \cite{Sch98}.
For $n\ge4$ they fill a locus of dimension $3n<n^{2}$. Koike proved the
analogous statement for a family of sixfolds for $\QQ(\sqrt{-1})$
\cite{Koi04}.

In Chapter~\ref{ch:remaining} we prove that cycles built from a curve with
an automorphism of order three in the obvious ways, through diagonals,
graphs of the automorphism and the group law, never represent a Weil class
(Proposition~\ref{prop:curvesbook}). Schoen's cycles therefore use something
beyond these constructions, and that is why they reach classes that divisors
cannot.

\section{Secant sheaves}

The most complete results are Markman's. He works with abelian varieties of
Weil type of the form $B\times\widehat B$, for an abelian variety $B$ of
dimension $n$ and its dual, with an action of $K$ that does not respect the
product. These are special members of their Weil families. On them he
constructs sheaves from the secant varieties of $B$ in a projective
embedding, the \emph{secant sheaves}, whose Chern characters have a nonzero
Weil component. He then shows that these sheaves deform along the whole
family, by an equivariant form of the semiregularity theorem of Bloch,
Buchweitz--Flenner and Pridham \cite{Blo72,BF03,Pri24}.

\begin{theorem}[Markman \cite{Mar25a,Mar25b}, preprints]\label{thm:markman}
\begin{enumerate}[label=(\roman*)]
\item The Weil classes are algebraic on every abelian sixfold of Weil type
whose hermitian form is split.
\item The Weil classes are algebraic on every abelian fourfold of Weil
type, for every imaginary quadratic field and every discriminant.
\end{enumerate}
\end{theorem}

Part (ii) follows from part (i) by Proposition~\ref{prop:lower}. Combined with
the analysis of Moonen and Zarhin \cite{MZ99}, it gives the Hodge
conjecture for every abelian variety of dimension at most five. At the time
of writing, Markman's papers are on the arXiv, and we mark the results as
preprint results accordingly.

It is worth seeing where secant sheaves sit with respect to the barrier of
Chapter~\ref{ch:remaining}. On $B\times\widehat B$ the endomorphism algebra is
much larger than $K$, the Lefschetz group is much smaller than $\U(V,H)$, and
the Weil classes are invariant under it. So nothing prevents a sheaf built
from line bundles from having a Weil component there. The difficulty is
moved into the deformation step, which is where all the work lies.

\section{Recent preprints}\label{sec:preprintsbook}

In September and October 2026, OpenAI released a series of preprints
announcing further cases \cite{OAI26d,OAI26e,OAI26g}. They have not been
refereed, and this book does not check them. Nothing outside this section
depends on them. The announced statements include:
\begin{enumerate}[label=(\alph*)]
\item the Hodge conjecture for every abelian variety of CM type
\cite{OAI26e};
\item the algebraicity of the Weil classes on every abelian eightfold of
Weil type whose hermitian form is split, for every imaginary quadratic
field \cite{OAI26d};
\item the Hodge conjecture for every power of an abelian sixfold of split
Weil type, and of an abelian variety of dimension at most five with an
action of an imaginary quadratic field \cite{OAI26g}.
\end{enumerate}
If (b) holds, Proposition~\ref{prop:lower} gives the Weil classes on every
abelian sixfold of Weil type \cite{BhH8}. Even then, the Weil classes in
dimension eight with non-split form, and in every dimension from ten on,
would remain open, and so would the classes on abelian varieties that are
not of Weil type and are not Lefschetz, which \cite{BhH8} lists.

\section{Summary}

\begin{center}
\begin{tabular}{@{}>{\raggedright\arraybackslash}p{3.6cm}>{\raggedright\arraybackslash}p{5.2cm}>{\raggedright\arraybackslash}p{5.0cm}@{}}
\toprule
Dimension $2n$ & Status of the Weil classes & Source\\
\midrule
$4$ & algebraic, all $K$ and discriminants & Markman (preprint); earlier cases Schoen\\
$6$, split form & algebraic & Markman (preprint)\\
$6$, $\QQ(\sqrt{-3})$, Prym family & algebraic & Schoen\\
$6$, one family for $\QQ(\sqrt{-1})$ & algebraic & Koike\\
$6$, other & open; announced via (b) & preprints\\
$8$, split form & open; announced & preprint (b)\\
$2n\ge8$, $\QQ(\sqrt{-3})$ & algebraic on a locus of dimension $3n$ & Schoen\\
$2n\ge8$, other & open, apart from special members & --\\
\bottomrule
\end{tabular}
\end{center}

In every row marked open, the Weil classes are algebraic on a dense set of
special members (Proposition~\ref{prop:densebook}), and the question is
whether that set has interior. That is the subject of the last chapter.
